% 编译：xelatex architecture.tex（需要 ctex 宏包，编译两遍以解析引用）
\documentclass[11pt,a4paper]{ctexart}
\usepackage[margin=2.5cm]{geometry}
\usepackage{amsmath,amssymb}
\usepackage{booktabs,tabularx}
\usepackage{enumitem}
\usepackage[hidelinks]{hyperref}
\usepackage{xurl}

\title{跨层误差传播模型：架构设计}
\date{2026年11月12日}

\begin{document}
\maketitle

\begin{abstract}
本文定义实现跨层误差传播模型的软件与硬件架构：四个误差接口 E1--E4 的数据格式，三个传播算子 $\mathcal P_{12}$、$\mathcal P_{23}$、$\mathcal P_{34}$ 的模块划分，各层误差源的注入点，每个算子的三种引擎（解析、仿真、硬件），以及组合与验证的数据流。目标、研究问题与数制约定见 \texttt{overview.tex}，公式见 \texttt{math-foundations.tex}。
\end{abstract}

\section{设计原则}
\begin{enumerate}[leftmargin=*]
  \item \textbf{以层间边界为中心}：模块按传播算子划分，而不是按层的内容划分。
  \item \textbf{接口显式且可替换}：每个接口的表示方式（信息量）是可配置的，以便研究接口充分性（R1）。
  \item \textbf{每个算子三种引擎}：解析引擎给出预测，仿真引擎与硬件引擎给出真值；三者使用相同的接口格式。
  \item \textbf{任意层可注入、任意层可观测}：支持“注入层 $\times$ 观测层”实验矩阵。
  \item \textbf{可追溯}：每个产物记录配置、代码版本与上游产物编号。
\end{enumerate}

\section{总体结构}

\begin{table}[h]
\centering
\small
\begin{tabularx}{\textwidth}{@{}lXl@{}}
\toprule
组成 & 职责 & 对应 \\
\midrule
误差源 $S_1$ & 器件层误差：位翻转注入（操作数、乘积、累加器、逻辑节点）；可选的时序违例 & E1 \\
算子 $\mathcal P_{12}$ & 位错误 $\to$ 单元输出误差：位权放大、操作数放大、逻辑掩蔽 & E1$\to$E2 \\
误差源 $S_2$ & 电路层误差：近似运算单元的设计误差 $e(a,b)$ & E2 \\
算子 $\mathcal P_{23}$ & 单元误差 $\to$ 网络层张量误差：累加、饱和、舍入、非线性、层间递推 & E2$\to$E3 \\
误差源 $S_3$ & 框架层误差：累加器截断、位宽缩减 & E3 \\
算子 $\mathcal P_{34}$ & 张量误差 $\to$ 任务质量：失配概率、大误差敏感度 & E3$\to$E4 \\
组合器 & 按配置串联算子，合并各层误差源 & E1$\to$E4 \\
验证器 & 注入层 $\times$ 观测层实验；计算充分性缺陷、组合误差、掩蔽与增益 & R1--R4 \\
\bottomrule
\end{tabularx}
\caption{总体结构}
\end{table}

\section{接口格式}
\label{sec:interfaces}
所有接口产物包含元数据头（产物编号、生成时间、git 提交号、配置、上游产物编号）；数据用 Parquet，元数据用 JSON。每个接口都有一个\textbf{表示级别}字段，取值决定携带的信息量。

\begin{table}[h]
\centering
\small
\begin{tabularx}{\textwidth}{@{}lXX@{}}
\toprule
接口 & 内容 & 表示级别（R1 的比较对象） \\
\midrule
E1 & 每个故障位置 $k$ 的出错概率 $\pi_k$（及置信区间）；位置类型（激活位、权重位、乘积位、累加器位、逻辑节点） & 单一总错误率 / 按位置 / 按位置且以输入跳变为条件 \\
E2 & 每个单元：$e(a,b)$ 查找表；故障引起的误差表 $\Delta_k(a,b)$；以权重为条件的矩 $\mu(a)$、$\sigma^2(a)$；大误差率 $\lambda(a)$ 与幅度分布 & 均匀分布下的矩 / 以权重为条件的矩 / 完整查找表 \\
E3 & 每个网络层：误差均值、协方差表示、大误差率 $\Lambda_l$；输出层 logits 误差的统计 & 对角方差 / 分组仿射算术 / 完整（仅仿真） \\
E4 & 精度下降、预测一致率、每个样本的翻转概率 & — \\
\bottomrule
\end{tabularx}
\caption{接口格式与表示级别}
\end{table}

\section{传播算子模块}

\subsection{$\mathcal P_{12}$：器件 $\to$ 电路}
\begin{table}[h]
\centering
\small
\begin{tabularx}{\textwidth}{@{}lX@{}}
\toprule
引擎 & 实现 \\
\midrule
解析 & 存储位翻转按位权与另一操作数解析计算 $\Delta_k(a,b)$；逻辑节点使用预先算好的掩蔽表 \\
仿真 & 门级网表仿真（Verilator 或 Icarus Verilog，基于 EvoApprox8b 的 Verilog）：逐节点注入、穷举 65536 种输入，得到 $\Delta_k(a,b)$ 与逻辑掩蔽概率 \\
硬件 & ZCU102 上的单元测试电路：注入掩码寄存器控制位翻转，读回输出 \\
\bottomrule
\end{tabularx}
\caption{$\mathcal P_{12}$ 的三种引擎}
\end{table}

\subsection{$\mathcal P_{23}$：电路 $\to$ 框架}
\begin{table}[h]
\centering
\small
\begin{tabularx}{\textwidth}{@{}lX@{}}
\toprule
引擎 & 实现 \\
\midrule
解析 & 逐层矩传播：点积均值方差、截断与饱和、舍入掩蔽、ReLU 闭式矩、BatchNorm 缩放、残差、池化；协方差表示可选对角或分组概率仿射算术；大误差记录 $\Lambda_l$ \\
仿真 & PyTorch 中用查找表替换乘法（AdaPT 或自写 CUDA 算子），并在操作数、乘积、累加器上挂故障注入钩子；逐比特复现硬件 \\
硬件 & ZCU102 上的 INT8 脉动阵列：乘法器以查找表（BRAM）实现，可整表替换；操作数、乘积、累加器寄存器上有注入掩码；ARM 端逐层调度并经 DMA 读回每层输出 \\
\bottomrule
\end{tabularx}
\caption{$\mathcal P_{23}$ 的三种引擎}
\end{table}

\subsection{$\mathcal P_{34}$：框架 $\to$ 应用}
\begin{table}[h]
\centering
\small
\begin{tabularx}{\textwidth}{@{}lX@{}}
\toprule
引擎 & 实现 \\
\midrule
解析 & 基于失配概率（含偏置与协方差）计算小误差翻转概率；用敏感度 $s_l$ 计算大误差翻转概率；合并为精度下降与一致率 \\
仿真 / 硬件 & 在实际输出上统计配对的翻转样本数，计算精度下降与一致率 \\
\bottomrule
\end{tabularx}
\caption{$\mathcal P_{34}$ 的引擎}
\end{table}

\paragraph{敏感度 $s_l$ 的估计.} 用仿真引擎在第 $l$ 层随机注入大误差，统计输出改变的比例；每个模型只做一次，结果入库。

\section{误差注入点}
\begin{table}[h]
\centering
\small
\begin{tabularx}{\textwidth}{@{}lXX@{}}
\toprule
误差源 & 注入位置 & 控制参数 \\
\midrule
$S_1$ 位翻转 & 激活、权重、乘积、累加器寄存器的指定位；网表内部节点（仅单元级） & 位置、每次运算出错概率、随机种子 \\
$S_1$ 时序违例（可选） & 超频被测时钟域 & 频率、电压、温度 \\
$S_2$ 近似单元 & 每个网络层选择一个乘法器查找表 & 每层乘法器编号 \\
$S_3$ 累加器 & 部分和截断低 $t$ 位；累加器位宽 $B$ & $t$、$B$ \\
\bottomrule
\end{tabularx}
\caption{误差注入点}
\end{table}

硬件上的位翻转注入用 LFSR 产生随机注入掩码，按设定概率与寄存器值异或；随机种子由 ARM 端设置，保证可复现。

\section{数据流}

\subsection{预测流}
给定配置（各层误差源参数、表示级别）：
\begin{enumerate}[leftmargin=*]
  \item 读取或计算 E1；
  \item $\mathcal P_{12}$ 得到故障引起的单元误差，与 $S_2$ 合并成 E2（在单元级合并，保留操作数依赖）；
  \item $\mathcal P_{23}$ 逐层传播，并在对应网络层加入 $S_3$，得到 E3；
  \item $\mathcal P_{34}$ 得到 E4。
\end{enumerate}

\subsection{真值流}
相同配置下，仿真引擎（逐比特）与硬件引擎分别运行，在 E2、E3（每个网络层）、E4 处导出实测值。仿真与硬件在查找表方案下应逐比特一致，作为正确性检查。

\subsection{验证流}
对每个“注入层 $\times$ 观测层”组合计算：
\begin{itemize}[leftmargin=*]
  \item 单个算子的充分性缺陷 $\eta_i$（以真值输入驱动单个算子，与真值输出比较）；
  \item 组合误差（以组合模型的输出与端到端真值比较）及其理论界；
  \item 各边界的掩蔽概率、能量增益、偏置传递。
\end{itemize}

\subsection{反馈流}
应用层在参考模型上统计各网络层的激活分布 $p_x$，提供给 E2 计算以权重为条件的矩。

\section{数制约定}
按总纲：权重有符号 INT8 对称零点 0；激活无符号 8 位零点 0；首末层精确；累加器 32 位为基准；重量化四舍五入；有符号乘法用符号—幅值包装无符号近似乘法器。

\section{代码仓库结构（建议）}
\begin{table}[h]
\centering
\small
\begin{tabularx}{\textwidth}{@{}lX@{}}
\toprule
目录 & 内容 \\
\midrule
\texttt{interfaces/} & E1--E4 的数据结构、表示级别、Parquet 读写、元数据 \\
\texttt{sources/} & $S_1$、$S_2$、$S_3$ 误差源的参数化与注入 \\
\texttt{operators/p12/} & 器件$\to$电路：解析、门级仿真 \\
\texttt{operators/p23/} & 电路$\to$框架：矩传播（对角、仿射算术）、查找表仿真 \\
\texttt{operators/p34/} & 框架$\to$应用：失配概率、敏感度估计 \\
\texttt{compose/} & 组合器 \\
\texttt{validate/} & 注入层 $\times$ 观测层实验、充分性缺陷、组合误差、掩蔽与增益 \\
\texttt{hw/} & HLS 与 RTL、注入掩码逻辑、Vivado 脚本、板上驱动 \\
\texttt{configs/}、\texttt{experiments/} & 配置与实验脚本 \\
\texttt{data/} & 测量数据与库（不纳入 git） \\
\bottomrule
\end{tabularx}
\caption{建议的代码仓库结构}
\end{table}

\end{document}
